% --------------------------------------------------------------------- %
%  CentralImmo -- Fichier de diagrammes TikZ autonomes pour le mémoire
%  Spécialement configuré pour compiler sur Overleaf en pdflatex
%  Couleurs thématiques : Vert Émeraude (#00695C) et Or Premium (#D4AF37)
% --------------------------------------------------------------------- %

\documentclass[11pt,a4paper]{report}

\usepackage[utf8]{inputenc}
\usepackage[T1]{fontenc}
\usepackage[french]{babel} % Gère les accents français nativement
\usepackage[margin=1.8cm,top=2cm,bottom=2cm]{geometry}
\usepackage{xcolor}
\usepackage{tikz}
\usepackage{caption}
\usepackage{float}
\usepackage{enumitem}

% Bibliothèques TikZ standard complètes
\usetikzlibrary{positioning, arrows.meta, shapes.geometric, fit, calc, shadows, decorations.pathreplacing, backgrounds}

% --- Palette de Couleurs CentralImmo ---
\definecolor{ciGreen}{HTML}{00695C}       % Vert Émeraude (Principal)
\definecolor{ciGreenLight}{HTML}{E0F2F1}  % Vert très clair pour les fonds
\definecolor{ciGold}{HTML}{D4AF37}        % Or Premium (Secondaire)
\definecolor{ciGoldLight}{HTML}{FCF8E3}   % Or très clair pour les alertes
\definecolor{ciDark}{HTML}{212121}        % Gris très foncé (presque noir) pour le texte
\definecolor{ciLightGray}{HTML}{F5F5F5}   % Gris clair pour les conteneurs généraux
\definecolor{ciBorder}{HTML}{757575}      % Gris pour les bordures secondaires
\definecolor{ciRed}{HTML}{C62828}         % Rouge discret pour le rejet
\definecolor{ciBlue}{HTML}{1976D2}        % Bleu discret pour la recherche

% Styles TikZ globaux réutilisables (Optimisés pour éviter les masquages et chevauchements)
\tikzset{
  % Noeuds de base
  base-node/.style={
    draw=ciDark, line width=0.8pt, rounded corners=4px,
    font=\small\sffamily, align=center, fill=white,
    inner sep=6pt, drop shadow={opacity=0.12, shadow xshift=1pt, shadow yshift=-1.2pt}
  },
  % Noeuds colorés par rôle
  node-green/.style={base-node, draw=ciGreen, fill=ciGreenLight, text=ciDark},
  node-gold/.style={base-node, draw=ciGold!80!black, fill=ciGoldLight, text=ciDark},
  node-white/.style={base-node, draw=ciBorder, fill=white, text=ciDark},
  node-red/.style={base-node, draw=ciRed, fill=ciRed!8, text=ciDark},
  node-blue/.style={base-node, draw=ciBlue, fill=ciBlue!10, text=ciDark},
  % Flèches de flux
  flow-arrow/.style={-Stealth, line width=0.9pt, draw=ciDark},
  flow-arrow-dashed/.style={-Stealth, line width=0.9pt, draw=ciGreen, dashed},
  flow-arrow-gold/.style={-Stealth, line width=1.1pt, draw=ciGold!80!black},
  % Textes sur les flèches
  label-style/.style={font=\scriptsize\sffamily, fill=white, inner sep=1.5pt, text=ciDark},
  % Conteneurs en fond (Remplissage retiré pour éviter de cacher le texte sous Overleaf)
  group-box/.style={
    draw=ciBorder, dotted, line width=0.8pt, rounded corners=6px,
    fill=none, inner sep=8pt
  },
  group-box-green/.style={
    draw=ciGreen, dashed, line width=0.8pt, rounded corners=6px,
    fill=none, inner sep=8pt
  }
}

\begin{document}

\title{\bfseries Recueil de Diagrammes Conceptuels et Techniques \\ \large Projet CentralImmo (Agrégateur Immobilier)}
\author{Bryan Azebaze}
\date{Soutenance de Licence Professionnelle -- Génie Logiciel}
\maketitle

\thispagestyle{empty}
\newpage

\chapter*{Introduction et Guide d'intégration}
Ce document propose **quinze diagrammes** représentatifs de la plateforme **CentralImmo** codés entièrement en TikZ. 
Tous les diagrammes ont été optimisés pour tenir dans la largeur d'une page A4 classique et éviter les chevauchements ou les boîtes opaques vides constatés lors de compilations LaTeX basiques.

\paragraph{Recommandation essentielle :} 
Pour que les fonds transparents et les regroupements s'affichent correctement, ajoutez bien \texttt{backgrounds} dans vos librairies TikZ de votre document principal :
\begin{verbatim}
\usetikzlibrary{positioning, arrows.meta, shapes.geometric, fit, calc, shadows, backgrounds}
\end{verbatim}

\newpage

% ================================================================== %
% DIAGRAMME 1 : ARCHITECTURE DU SYSTEME (CORRIGÉ & LINÉAIRE)
% ================================================================== %
\section*{Diagramme 1 : Architecture globale du système (Flux vertical)}
\label{diag:architecture}
Ce schéma montre le flux logique vertical : de la collecte à la diffusion finale.

\begin{figure}[H]
\centering
\begin{tikzpicture}[node distance=8mm and 8mm, scale=0.95, every node/.style={scale=0.95}]

  % --- Plan Collecte (Scrapers) ---
  \node[node-gold] (mapiole) {Scraper Dédié\\Mapiole (Python)};
  \node[node-gold, right=8mm of mapiole] (kasastay) {Scraper Dédié\\Kasastay (Python)};
  \node[node-gold, right=8mm of kasastay] (universal) {Scraper Universel\\(Heuristiques)};

  % --- Plan Application (API - FastAPI) ---
  \node[node-green, below=15mm of kasastay, text width=6cm] (fastapi) {
    \textbf{Serveur Backend (FastAPI)}\\
    \footnotesize
    $\bullet$ Ingestion \& déduplication (DCS)\\
    $\bullet$ API de recherche naturelle \& isochrone
  };

  % --- Plan Persistance (Base de données) ---
  \node[node-white, right=12mm of fastapi, text width=4.5cm] (db) {
    \textbf{Base PostgreSQL}\\
    \color{ciDark}\scriptsize
    $\bullet$ \texttt{raw\_listings}\\
    $\bullet$ \texttt{canonical\_properties}\\
    $\bullet$ \texttt{listing\_history}\\
    $\bullet$ \texttt{locations}
  };

  % --- Plan Présentation (Client Mobile Flutter) ---
  \node[node-blue, below=15mm of fastapi, text width=5cm] (client) {
    \textbf{Application Mobile (Flutter)}\\
    \footnotesize Interface unifiée sans logo de source en page d'accueil (Modèle Trivago)
  };

  % --- Groupes (Fonds transparents dessinés sur le calque de fond) ---
  \begin{scope}[on background layer]
    \node[group-box, fit=(mapiole) (kasastay) (universal), label={[font=\footnotesize\bfseries\sffamily, text=ciGold!80!black]above:1. PHASE DE COLLECTE ET EXTRACTION}] (box-collecte) {};
  \end{scope}

  % --- Liaisons ---
  \draw[flow-arrow] (mapiole.south) -- ++(0,-4mm) -| ([xshift=-1.5cm]fastapi.north);
  \draw[flow-arrow] (kasastay.south) -- (fastapi.north);
  \draw[flow-arrow] (universal.south) -- ++(0,-4mm) -| ([xshift=1.5cm]fastapi.north);

  \draw[flow-arrow, draw=ciGreen, line width=1.1pt] (fastapi) -- node[label-style, above] {SQL} (db);
  \draw[flow-arrow, draw=ciBlue, line width=1.1pt] (client) -- node[label-style, right] {HTTP Get / Search} (fastapi);
  \draw[flow-arrow-dashed, draw=ciBlue] (fastapi.south west) -- node[label-style, left] {JSON Response} (client.north west);

\end{tikzpicture}
\caption{Architecture globale de récupération, de stockage et de distribution.}
\end{figure}

\newpage

% ================================================================== %
% DIAGRAMME 2 : CAS D'UTILISATION (RÉSOLU - SANS FOND BLANC MASQUANT)
% ================================================================== %
\section*{Diagramme 2 : Cas d'utilisation global et Acteurs}
\label{diag:usecase}
Représentation simplifiée des interactions des acteurs clés avec la plateforme.

\begin{figure}[H]
\centering
\begin{tikzpicture}[node distance=8mm and 15mm, scale=0.9, every node/.style={scale=0.9}]

  % --- Acteurs ---
  \node[align=center] (user) {
    \begin{tikzpicture}[scale=0.5]
      \draw[fill=white, line width=0.8pt] (0,1) circle (0.25);
      \draw[line width=0.8pt] (0,0.75) -- (0,0);
      \draw[line width=0.8pt] (-0.35,0.45) -- (0.35,0.45);
      \draw[line width=0.8pt] (0,0) -- (-0.25,-0.5);
      \draw[line width=0.8pt] (0,0) -- (0.25,-0.5);
    \end{tikzpicture}\\
    \small Utilisateur Final\\(Chercheur de logement)
  };

  \node[align=center, right=8cm of user] (admin) {
    \begin{tikzpicture}[scale=0.5]
      \draw[fill=white, line width=0.8pt] (0,1) circle (0.25);
      \draw[line width=0.8pt] (0,0.75) -- (0,0);
      \draw[line width=0.8pt] (-0.35,0.45) -- (0.35,0.45);
      \draw[line width=0.8pt] (0,0) -- (-0.25,-0.5);
      \draw[line width=0.8pt] (0,0) -- (0.25,-0.5);
    \end{tikzpicture}\\
    \small Administrateur SSH\\(Modérateur technique)
  };

  % --- Cas d'utilisation ---
  \node[node-green, right=1.6cm of user, yshift=2.2cm, text width=3.5cm] (uc1) {Rechercher par critères};
  \node[node-green, right=1.6cm of user, yshift=0.7cm, text width=3.5cm] (uc2) {Comparer les offres (Trivago)};
  \node[node-green, right=1.6cm of user, yshift=-0.7cm, text width=3.5cm] (uc3) {Recherche par temps de trajet};
  \node[node-green, right=1.6cm of user, yshift=-2.2cm, text width=3.5cm] (uc4) {Consulter les tendances prix};

  \node[node-gold, left=1.6cm of admin, yshift=1cm, text width=3.5cm] (uc5) {Valider/Modérer les annonces};
  \node[node-gold, left=1.6cm of admin, yshift=-1cm, text width=3.5cm] (uc6) {Superviser les scrapers};

  % Frontière de l'Application (fill=none pour ne pas masquer les ellipses)
  \begin{scope}[on background layer]
    \node[group-box, fit=(uc1) (uc2) (uc3) (uc4) (uc5) (uc6), label={[font=\small\bfseries\sffamily, text=ciGreen]above:SERVICES DE L'APPLICATION CENTRALIMMO}] (system-boundary) {};
  \end{scope}

  % Connecteurs
  \draw[flow-arrow] (user) -- (uc1);
  \draw[flow-arrow] (user) -- (uc2);
  \draw[flow-arrow] (user) -- (uc3);
  \draw[flow-arrow] (user) -- (uc4);

  \draw[flow-arrow-gold] (admin) -- (uc5);
  \draw[flow-arrow-gold] (admin) -- (uc6);

\end{tikzpicture}
\caption{Diagramme de cas d'utilisation (Limite du système en pointillés transparents).}
\end{figure}

\newpage

% ================================================================== %
% DIAGRAMME 3 : MODELE TRIVAGO STYLE (EDITIONS ET COMPRESSIONS)
% ================================================================== %
\section*{Diagramme 3 : Modèle de Déduplication Comparatif (Modèle Trivago)}
\label{diag:trivago}
Ce schéma illustre la fusion de plusieurs annonces vers une unique super-annonce référente.

\begin{figure}[H]
\centering
\begin{tikzpicture}[node distance=8mm, scale=0.9, every node/.style={scale=0.9}]

  % --- Annonces Sources Originales ---
  \node[node-white, text width=4cm] (src1) {
    \textbf{Kasastay (Source 1)}\\
    \scriptsize
    $\bullet$ Titre: "Studio meublé Akwa"\\
    $\bullet$ Prix: \textbf{45 000 FCFA/nuit}
  };

  \node[node-white, below=4mm of src1, text width=4cm] (src2) {
    \textbf{Mapiole (Source 2)}\\
    \scriptsize
    $\bullet$ Titre: "Studio Moderne Akwa"\\
    $\bullet$ Prix: \textbf{50 000 FCFA/nuit}
  };

  \node[node-white, below=4mm of src2, text width=4cm] (src3) {
    \textbf{Universal (Source 3)}\\
    \scriptsize
    $\bullet$ Titre: "Superbe Studio Akwa"\\
    $\bullet$ Prix: \textbf{43 000 FCFA/nuit}
  };

  % --- Étape Déduplication / Ingestion (FastAPI) ---
  \node[node-gold, right=1.6cm of src2, text width=3cm] (fusion) {
    \textbf{Pipeline de Fusion}\\
    \footnotesize
    Calcul DCS $\approx$ 91\%\\
    (Lieu et Titre similaires)
  };

  % --- Super Annonce Canonique ---
  \node[node-green, right=1.6cm of fusion, text width=5cm] (canonical) {
    \textbf{SUPER-ANNONCE (Unique)}\\
    \small
    \textbf{Studio meublé moderne à Akwa}\\
    \color{ciGreen}\textbf{Meilleur Prix : 43 000 FCFA / nuit}\\
    \color{ciDark}\scriptsize
    $\bullet$ Photos agrégées sans doublon\\
    \color{ciBorder}\rule{\linewidth}{0.3pt}\\
    \color{ciDark}\textbf{Autres Offres liées :}\\
    - Kasastay : 45 000 FCFA\\
    - Mapiole : 50 000 FCFA
  };

  % --- Connecteurs ---
  \draw[flow-arrow] (src1.east) -- (fusion.west);
  \draw[flow-arrow] (src2.east) -- (fusion.west);
  \draw[flow-arrow] (src3.east) -- (fusion.west);

  \draw[flow-arrow, draw=ciGreen, line width=1.3pt] (fusion.east) -- node[label-style, above] {Génère} (canonical.west);

  % Représentation visuelle des liens
  \draw[flow-arrow-dashed, draw=ciGreen, bend right=20] (canonical.south west) to  (src1.south east);
  \draw[flow-arrow-dashed, draw=ciGreen] (canonical.west) to (src2.east);
  \draw[flow-arrow-dashed, draw=ciGreen, bend left=20] (canonical.north west) to (src3.north east);

\end{tikzpicture}
\caption{Processus Trivago de regroupement d'offres immobilières équivalentes.}
\end{figure}

\newpage

% ================================================================== %
% DIAGRAMME 4 : DCS WEIGHTS / PONDÉRATIONS
% ================================================================== %
\section*{Diagramme 4 : Algorithme DCS (Duplicate Confidence Score)}
\label{diag:dcs}
Pondérations intégrées dans l'évaluation d'un doublon entre un brouillon et une annonce référente.

\begin{figure}[H]
\centering
\begin{tikzpicture}[scale=0.9, every node/.style={scale=0.9}]

  % --- Facteurs de Scoring ---
  \node[node-white, font=\scriptsize\sffamily, text width=4.2cm] (f1)  {Similitude du Titre (Token ratio)};
  \node[node-white, font=\scriptsize\sffamily, text width=4.2cm, below=3mm of f1] (f2)  {Proximité du Prix ($\pm15$\%)};
  \node[node-white, font=\scriptsize\sffamily, text width=4.2cm, below=3mm of f2] (f3)  {Type de bien (Villa, Studio...)};
  \node[node-white, font=\scriptsize\sffamily, text width=4.2cm, below=3mm of f3] (f4)  {Surface en m² (Tolérance 15\%)};
  \node[node-white, font=\scriptsize\sffamily, text width=4.2cm, below=3mm of f4] (f5)  {Nb Chambres (Comparaison exacte)};
  \node[node-white, font=\scriptsize\sffamily, text width=4.2cm, below=3mm of f5] (f6)  {Géolocalisation résolue};
  \node[node-white, font=\scriptsize\sffamily, text width=4.2cm, below=3mm of f6] (f7)  {Image overlap (MD5)};

  % --- Poids des coefficients ---
  \node[node-gold, right=1.5cm of f1] (p1) {Poids : 30\%};
  \node[node-gold, right=1.5cm of f2] (p2) {Poids : 20\%};
  \node[node-gold, right=1.5cm of f3] (p3) {Poids : 15\%};
  \node[node-gold, right=1.5cm of f4] (p4) {Poids : 15\%};
  \node[node-gold, right=1.5cm of f5] (p5) {Poids : 10\%};
  \node[node-gold, right=1.5cm of f6] (p6) {Poids : 5\%};
  \node[node-gold, right=1.5cm of f7] (p7) {Poids : 5\%};

  % --- Bloc Sommateur ---
  \node[node-green, shape=circle, right=1.5cm of p4, inner sep=4pt] (sum) {\textbf{Somme \\ Pondérée}};

  % --- Verdict final ---
  \node[node-green, right=1.5cm of sum, xshift=10pt, yshift=1.2cm, text width=2.6cm] (conf) {\textbf{Match Confirmé}\\Score $\ge 85\%$};
  \node[node-gold, right=1.5cm of sum, xshift=10pt, text width=2.6cm] (prob) {\textbf{Match Probable}\\ $65\% \le$ Score $< 85\%$};
  \node[node-red, right=1.5cm of sum, xshift=10pt, yshift=-1.2cm, text width=2.6cm] (newp) {\textbf{Nouvelle Fiche}\\Score $< 65\%$};

  % --- Liaisons ---
  \foreach \i in {1,...,7} {
    \draw[flow-arrow] (f\i) -- (p\i);
    \draw[flow-arrow, draw=ciGold] (p\i.east) -- (sum.west);
  }
  
  \draw[flow-arrow, draw=ciGreen] (sum.east) -- (conf.west);
  \draw[flow-arrow, draw=ciGold] (sum.east) -- (prob.west);
  \draw[flow-arrow, draw=ciRed] (sum.east) -- (newp.west);

\end{tikzpicture}
\caption{Algorithme de calcul du coefficient de similitude DCS.}
\end{figure}

\newpage

% ================================================================== %
% DIAGRAMME 5 : STRUCTURE DE DONNEES (COMPRIS PAR LE JURY)
% ================================================================== %
\section*{Diagramme 5 : Structure des données SQL (MCD Simplifié)}
\label{diag:mcd}
Ce diagramme montre comment la base PostgreSQL stocke l'historique d'une annonce.

\begin{figure}[H]
\centering
\begin{tikzpicture}[node distance=1.2cm and 1.8cm, scale=0.85, every node/.style={scale=0.85}]

  % Entité Source
  \node[node-gold, text width=4.2cm] (src) {
    \textbf{Entity: Sources}\\
    \color{ciGreen}\scriptsize PK: id (integer)\\
    \color{ciDark}\tiny
    $\bullet$ slug (unique string)\\
    $\bullet$ display\_name (string)\\
    $\bullet$ adapter\_kind (string)
  };

  % Entité RawListing
  \node[node-green, below=of src, text width=4.2cm] (raw) {
    \textbf{Entity: RawListings}\\
    \color{ciGreen}\scriptsize PK: id (bigint)\\
    \color{ciDark}\tiny
    $\bullet$ url\_source (text)\\
    $\bullet$ url\_hash (unique)\\
    $\bullet$ title\_raw (string)\\
    $\bullet$ price\_parsed (bigint)\\
    \color{ciBlue}FK: source\_id\\
    \color{ciBlue}FK: canonical\_property\_id
  };

  % Entité CanonicalProperty
  \node[node-green, right=of raw, text width=4.5cm] (canon) {
    \textbf{Entity: CanonicalProperties}\\
    \color{ciGreen}\scriptsize PK: id (bigint)\\
    \color{ciDark}\tiny
    $\bullet$ match\_key (string)\\
    $\bullet$ title\_canonical (string)\\
    $\bullet$ current\_best\_price (bigint)\\
    \color{ciBlue}FK: location\_id
  };

  % Entité Location
  \node[node-white, above=of canon, text width=4.2cm] (loc) {
    \textbf{Entity: Locations}\\
    \color{ciGreen}\scriptsize PK: id (bigint)\\
    \color{ciDark}\tiny
    $\bullet$ city (string)\\
    $\bullet$ neighborhood (string)\\
    $\bullet$ slug (string)
  };

  % Entité ListingHistory
  \node[node-red, right=of canon, text width=4.2cm] (hist) {
    \textbf{Entity: ListingHistory}\\
    \color{ciGreen}\scriptsize PK: id (bigint)\\
    \color{ciDark}\tiny
    $\bullet$ event\_type (string)\\
    $\bullet$ price\_observed (bigint)\\
    $\bullet$ observed\_at (datetime)
  };

  % Relations
  \draw[flow-arrow, draw=ciGold] (src.south) -- node[label-style, left] {1..N} (raw.north);
  \draw[flow-arrow, draw=ciGreen] (raw.east) -- node[label-style, above] {N..0-1} (canon.west);
  \draw[flow-arrow, draw=ciGreen] (loc.south) -- node[label-style, right] {1..N} (canon.north);
  \draw[flow-arrow, draw=ciRed] (canon.east) -- node[label-style, above] {1..N} (hist.west);
  \draw[flow-arrow, draw=ciRed, bend right=20] (raw.south) -- ++(0,-0.5cm) -| node[label-style, pos=0.3, below] {1..N} (hist.south);

\end{tikzpicture}
\caption{Modèle Conceptuel de Données (MCD).}
\end{figure}

\newpage

% ================================================================== %
% DIAGRAMME 6 : DIAGRAMME DE SEQUENCE DE L'INGESTION
% ================================================================== %
\section*{Diagramme 6 : Séquence d'extraction et intégration des données}
\label{diag:sequence}
Cycle de vie temporel d'un crawl d'annonces jusqu'à leur mise à disposition.

\begin{figure}[H]
\centering
\begin{tikzpicture}[scale=0.85, every node/.style={scale=0.85}]

  % --- Acteurs du flux ---
  \draw[line width=1pt] (0,0) -- (0,-7.5);
  \node[node-gold] at (0,0.3) {Sites Web};

  \draw[line width=1pt] (3.2,0) -- (3.2,-7.5);
  \node[node-gold] at (3.2,0.3) {Scrapers};

  \draw[line width=1pt] (7.2,0) -- (7.2,-7.5);
  \node[node-green] at (7.2,0.3) {FastAPI Ingest};

  \draw[line width=1pt] (11.2,0) -- (11.2,-7.5);
  \node[node-white] at (11.2,0.3) {PostgreSQL};

  \draw[line width=1pt] (14.2,0) -- (14.2,-7.5);
  \node[node-blue] at (14.2,0.3) {App Flutter};

  % --- Les Messages ---
  \draw[flow-arrow] (3.2,-1) -- node[label-style, above] {Requêtes HTTP (Jitter)} (0,-1);
  \draw[flow-arrow, dashed] (0,-1.7) -- node[label-style, above] {HTML brut} (3.2,-1.7);

  \draw[flow-arrow] (3.2,-2.4) -- node[label-style, above] {RawListingDraft} (7.2,-2.4);
  \draw[flow-arrow] (7.2,-3.2) -- node[label-style, above] {Rechercher clé unique} (11.2,-3.2);
  \draw[flow-arrow, dashed] (11.2,-3.8) -- node[label-style, above] {Retour clé existante / vide} (7.2,-3.8);

  \draw[flow-arrow] (7.2,-4.6) -- node[label-style, above] {Calcul DCS et liaison} (11.2,-4.6);
  \draw[flow-arrow] (7.2,-5.4) -- node[label-style, above] {INSERT/UPDATE Historique} (11.2,-5.4);

  \draw[flow-arrow, draw=ciBlue] (14.2,-6.2) -- node[label-style, above] {Appel GET /annonces} (7.2,-6.2);
  \draw[flow-arrow, draw=ciBlue, dashed] (7.2,-6.8) -- node[label-style, above] {Réponse JSON (Dédupliquée)} (14.2,-6.8);

\end{tikzpicture}
\caption{Flux de synchronisation asynchrone des données.}
\end{figure}

\newpage

% ================================================================== %
% DIAGRAMME 7 : DECISION D'INGESTION / COMPRESSE & COMPACTE
% ================================================================== %
\section*{Diagramme 7 : Organigramme d'Ingestion et Idempotence}
\label{diag:organigramme}
Optimisation du flux décisionnel pour éviter les écritures en doublon (cette version est compactée pour tenir sur une seule page).

\begin{figure}[H]
\centering
\begin{tikzpicture}[node distance=8mm, scale=0.88, every node/.style={scale=0.88}]

  % --- Points d'entrée ---
  \node[node-gold] (start) {Donnée extraite brute\\(\texttt{RawListingDraft})};
  \node[base-node, draw=ciDark, shape=diamond, aspect=2.2, below=6mm of start] (dec1) {L'URL existe en base ?};

  % --- Chemin RE-CRAWL (Gauche) ---
  \node[node-white, left=1.2cm of dec1, text width=3.8cm] (recrawl) {\textbf{Déjà vu}\\Mise à jour \texttt{last\_seen\_at}};
  \node[base-node, draw=ciDark, shape=diamond, aspect=2.2, below=6mm of recrawl] (dec2) {Le prix a changé ?};
  \node[node-green, below=6mm of dec2, text width=3.8cm] (hist-price) {Ajout événement\\dans \texttt{listing\_history}};

  % --- Chemin NOUVELLE ANNONCE (Droite) ---
  \node[node-white, right=1.2cm of dec1, text width=3.8cm] (new) {\textbf{Inconnu}\\Résolution localisation};
  \node[node-green, below=8mm of new, text width=3.8cm] (match) {Calcul DCS du doublon};
  \node[node-blue, below=8mm of match, text width=3.8cm] (insert) {Insertion \& Liaison parent};

  % --- Connecteurs ---
  \draw[flow-arrow] (start) -- (dec1);
  \draw[flow-arrow] (dec1.west) -- node[label-style, above] {Oui} (recrawl.east);
  \draw[flow-arrow] (dec1.east) -- node[label-style, above] {Non} (new.west);

  \draw[flow-arrow] (recrawl) -- (dec2);
  \draw[flow-arrow] (dec2.south) -- node[label-style, right] {Oui} (hist-price.north);
  \draw[flow-arrow] (dec2.west) -- node[label-style, above] {Non} ++(-1.2cm, 0) |- (hist-price.west);

  \draw[flow-arrow] (new) -- (match);
  \draw[flow-arrow] (match) -- (insert);

\end{tikzpicture}
\caption{Logique décisionnelle assurant l'idempotence des enregistrements.}
\end{figure}

\newpage

% ================================================================== %
% DIAGRAMME 8 : RECHERCHE IA CONVERSATIONNELLE
% ================================================================== %
\section*{Diagramme 8 : Fonctionnement de la Recherche IA Intégrée}
\label{diag:searchia}
Comment une requête rédigée en langage naturel est traduite en filtres de base de données.

\begin{figure}[H]
\centering
\begin{tikzpicture}[node distance=8mm]

  \node[node-blue] (query) {Saisie libre utilisateur :\\ \emph{"Appartement à Douala Akwa entre 30 et 50 millions"}};
  
  \node[node-white, below=6mm of query, text width=6cm] (api) {FastAPI Router (\texttt{api/search.py})\\Analyse de la clef DashScope};

  % AI
  \node[node-green, below left=10mm and 2mm of api, text width=4.5cm] (ai) {
    \textbf{Traitement IA (Qwen)}\\
    Extraction structurée :\\
    \footnotesize\color{ciGreen}
    \texttt{\{type: "Appartement",\\ price\_min: 30000000,\\ price\_max: 50000000,\\ city: "Douala",\\ neighborhood: "Akwa"\}}
  };

  % Regex
  \node[node-red, below right=10mm and 2mm of api, text width=4.5cm] (regex) {
    \textbf{Fallback Expressions régulières}\\
    Décomposition par mots-clés :\\
    \footnotesize\color{ciRed}
    Extraction basique des chiffres,\\ villes et types de biens
  };

  % Query SQL
  \node[node-white, below=4.2cm of api, text width=7.2cm] (sql) {
    \textbf{Exécution ORM (SQLAlchemy)}\\
    Génération et exécution de la requête SQL correspondante
  };

  \node[node-blue, below=6mm of sql] (results) {Affichage des fiches correspondantes};

  % --- Connecteurs ---
  \draw[flow-arrow] (query) -- (api);
  \draw[flow-arrow] (api.south west) -- node[label-style, left] {Jetons valides} (ai.north);
  \draw[flow-arrow] (api.south east) -- node[label-style, right] {Mode hors-ligne} (regex.north);
  
  \draw[flow-arrow, draw=ciGreen] (ai.south) -- ++(0,-4mm) -| ([xshift=-1.5cm]sql.north);
  \draw[flow-arrow, draw=ciRed] (regex.south) -- ++(0,-4mm) -| ([xshift=1.5cm]sql.north);
  \draw[flow-arrow] (sql) -- (results);

\end{tikzpicture}
\caption{Fonctionnement de la recherche naturelle intelligente.}
\end{figure}

\newpage

% ================================================================== %
% DIAGRAMME 9 : ISOCHRONE PIPELINE
% ================================================================== %
\section*{Diagramme 9 : Pipeline du calcul de déplacement par temps de trajet (Isochrone)}
\label{diag:isochrone}
Processus gérant la conversion d'un objectif de transport en résultats géolocalisés.

\begin{figure}[H]
\centering
\begin{tikzpicture}[node distance=8mm]

  \node[node-blue, text width=6cm] (user) {Saisie du point d'intérêt par l'utilisateur :\\ \emph{"Mon bureau" + "15 min de trajet"}};

  \node[node-white, below=5mm of user, text width=6cm] (nominatim) {
    \textbf{Géocodage (Nominatim API)}\\
    Résolution des coordonnées :\\
    $\to$ Lat / Lng du point de départ
  };

  \node[node-green, below=5mm of nominatim, text width=6cm] (osrm) {
    \textbf{Calculateur d'itinéraire (OSRM API)}\\
    Calcul de l'isochrone routier :\\
    $\to$ Polygone limite GeoJSON
  };

  \node[node-white, below=5mm of osrm, text width=6cm] (geom) {Détermination de la surface limite};

  \node[node-white, right=1.6cm of osrm, text width=4.5cm] (db) {
    \textbf{Requête SQL Géographique}\\
    Filtre les coordonnées des annonces à l'intérieur du polygone
  };

  \node[node-blue, below=12mm of geom, text width=5cm] (map) {Affichage des logements sur la carte de l'application (Flutter Map)};

  % --- Connecteurs ---
  \draw[flow-arrow] (user) -- (nominatim);
  \draw[flow-arrow] (nominatim) -- (osrm);
  \draw[flow-arrow] (osrm) -- (geom);
  \draw[flow-arrow, draw=ciGreen] (geom.east) -- (db.west);
  \draw[flow-arrow] (db.south) -- ++(0,-8mm) -| (map.east);
  \draw[flow-arrow] (geom.south) -- (map.north);

\end{tikzpicture}
\caption{Génération logicielle de l'affichage isochrone du territoire.}
\end{figure}

\newpage

% ================================================================== %
% DIAGRAMME 10 : SCREENFLOW NAVIGATION FLUTTER
% ================================================================== %
\section*{Diagramme 10 : Parcours utilisateur dans l'application mobile (Screenflow)}
\label{diag:screenflow}
Architecture globale de navigation ergonomique au sein de l'application cliente.

\begin{figure}[H]
\centering
\begin{tikzpicture}[node distance=6mm and 8mm, scale=0.88, every node/.style={scale=0.88}]

  \node[node-white, text width=3.2cm, minimum height=1.8cm] (onboarding) {
    \textbf{Écran d'Onboarding}\\
    \scriptsize Choisir la ville
  };

  \node[node-blue, right=of onboarding, text width=3.8cm, minimum height=2cm] (home) {
    \textbf{Écran d'Accueil}\\
    \scriptsize
    $\bullet$ Liste dédupliquée\\
    $\bullet$ \color{ciRed}Sans logos des sources\\
    $\bullet$ \color{ciDark}Bouton assistant IA
  };

  \node[node-green, right=of home, text width=4.2cm, minimum height=2.2cm] (detail) {
    \textbf{Détail de Logement}\\
    \scriptsize
    $\bullet$ Carousel Photos\\
    $\bullet$ \color{ciGreen}Meilleur tarif en gros\\
    $\bullet$ \color{ciDark}Bouton "Voir l'offre"\\
    $\bullet$ Liste comparative Trivago
  };

  \node[node-white, below=12mm of home, text width=3.8cm, minimum height=1.8cm] (map) {
    \textbf{Sélection Carte (Isochrone)}\\
    \scriptsize Filtre spatial sur la carte
  };

  % --- Liaisons ---
  \draw[flow-arrow] (onboarding) -- (home);
  \draw[flow-arrow, draw=ciGreen] (home.east) -- node[label-style, above] {Clic fiche} (detail.west);
  \draw[flow-arrow, draw=ciGreen, bend right=15] (detail.south west) to node[label-style, below] {Retour} (home.south east);

  \draw[flow-arrow] (home.south) -- (map.north);
  \draw[flow-arrow] (map.north west) -- (home.south west);
  \draw[flow-arrow, draw=ciGreen] (map.east) -- ++(5mm,0) |- (detail.south);

\end{tikzpicture}
\caption{Screenflow fonctionnel de l'application mobile.}
\end{figure}

\newpage

% ================================================================== %
% DIAGRAMME 11 : DEPLOIEMENT PHYSIQUE / CLOUD AWS (NOUVEAU)
% ================================================================== %
\section*{Diagramme 11 : Diagramme de Déploiement Cloud (Infrastructure de Production)}
\label{diag:deploiement}
Ce diagramme détaille comment la solution logicielle est déployée sur le Cloud d'AWS.

\begin{figure}[H]
\centering
\begin{tikzpicture}[node distance=8mm, scale=0.88, every node/.style={scale=0.88}]

  % --- Utilisateur ---
  \node[node-blue, text width=4cm] (user) {Utilisateur Final\\(Navigateur / Smartphone)};

  % --- Cloud Amazon Web Services ---
  \node[node-green, below=1.2cm of user, text width=12cm, minimum height=4.2cm, align=left] (vpc) {
    \quad \textbf{VPC Amazon Web Services (10.0.0.0/16)}\\[6pt]
    \qquad \textbf{Sous-réseau Public :}\\
    \qquad \quad $\bullet$ \texttt{Application Load Balancer} (ALB) - Redirection de trafic HTTPS \\[12pt]
    \qquad \textbf{Sous-réseau Privé :}\\
    \qquad \quad $\bullet$ \texttt{AWS EC2 Instance (t3.medium)}\\
    \qquad \quad \quad \footnotesize Linux + Uvicorn/FastAPI + Scrapers \quad \textbf{(Port 8000)} \\[10pt]
    \qquad \quad $\bullet$ \texttt{Amazon RDS PostgreSQL 16} (Base managée) \quad \textbf{(Port 5432)}
  };

  % Liaisons
  \draw[flow-arrow, line width=1pt] (user.south) -- node[label-style, right] {Port 443} ([yshift=3.2cm]vpc.south);

\end{tikzpicture}
\caption{Déploiement architectural sur les services cloud d'Amazon Web Services (AWS).}
\end{figure}

\newpage

% ================================================================== %
% DIAGRAMME 12 : ARBRE DE DECISION CONFIANCE UNIVERSAL SCRAPER (NOUVEAU)
% ================================================================== %
\section*{Diagramme 12 : Arbre de décision - Indice de confiance d'extraction}
\label{diag:confiance}
Comment le scraper évalue s'il doit automatiquement publier les annonces ou demander validation.

\begin{figure}[H]
\centering
\begin{tikzpicture}[node distance=8mm, scale=0.88, every node/.style={scale=0.88}]

  \node[node-gold] (start) {Extraction brute d'un site non certifié};
  
  \node[base-node, shape=diamond, aspect=2.2, below=6mm of start] (dec1) {Présence de coordonnées GPS ?};
  \node[base-node, shape=diamond, aspect=2.2, below=6mm of dec1] (dec2) {Le prix est-il cohérent?};

  \node[node-red, right=1.8cm of dec1, text width=3cm] (rec1) {Confiance faible\\(Gating: Rejet)};
  \node[node-white, left=1.8cm of dec2, text width=3.5cm] (rec2) {Confiance moyenne\\(Gating: \textbf{Pending} admin)};
  
  \node[node-green, below=8mm of dec2, text width=4cm] (auto) {Confiance élevée $\ge 0.75$\\(Gating: \textbf{Auto-Promoted})};

  % Connecteurs
  \draw[flow-arrow] (start) -- (dec1);
  \draw[flow-arrow] (dec1.east) -- node[label-style, above] {Non} (rec1.west);
  \draw[flow-arrow] (dec1.south) -- node[label-style, right] {Oui} (dec2.north);

  \draw[flow-arrow] (dec2.west) -- node[label-style, above] {Non} (rec2.east);
  \draw[flow-arrow] (dec2.south) -- node[label-style, right] {Oui} (auto.north);

\end{tikzpicture}
\caption{Calcul dynamique du score de confiance d'extraction pour le scraper universel.}
\end{figure}

\newpage

% ================================================================== %
% DIAGRAMME 13 : SCHEDULER BACKEND (NOUVEAU)
% ================================================================== %
\section*{Diagramme 13 : Planification des Tâches Asynchrones (APScheduler)}
\label{diag:scheduler}
Explication du cycle des crons gérés en arrière-plan du serveur FastAPI.

\begin{figure}[H]
\centering
\begin{tikzpicture}[node distance=8mm, scale=0.9, every node/.style={scale=0.9}]

  % Le démon APScheduler
  \node[node-green, text width=7.2cm] (daemon) {
    \textbf{APScheduler Daemon}\\
    \small S'exécute en arrière-plan du service FastAPI
  };

  % Tâches planifiées
  \node[node-white, below left=1.2cm and 2mm of daemon, text width=4.5cm] (job1) {
    \textbf{Crawl des Sources}\\
    \scriptsize
    $\bullet$ Planification : Tous les jours à 02h00\\
    $\bullet$ Action : Appel des scrapers
  };

  \node[node-white, below=1.2cm of daemon, text width=4.5cm] (job2) {
    \textbf{Nettoyage des Disparitions}\\
    \scriptsize
    $\bullet$ Planification : Toutes les 6 heures\\
    $\bullet$ Action : Marquer les supprimés
  };

  \node[node-white, below right=1.2cm and 2mm of daemon, text width=4.5cm] (job3) {
    \textbf{Indicateurs (Analytics)}\\
    \scriptsize
    $\bullet$ Planification : Chaque lundi à 03h00\\
    $\bullet$ Action : Recompute all scores
  };

  % Liaisons
  \draw[flow-arrow, draw=ciGreen] (daemon.south) -- ++(0,-6mm) -| (job1.north);
  \draw[flow-arrow, draw=ciGreen] (daemon.south) -- (job2.north);
  \draw[flow-arrow, draw=ciGreen] (daemon.south) -- ++(0,-6mm) -| (job3.north);

\end{tikzpicture}
\caption{Orchestration des tâches planifiées d'administration des données.}
\end{figure}

\newpage

% ================================================================== %
% DIAGRAMME 14 : PRICE ANALYZER (NOUVEAU)
% ================================================================== %
\section*{Diagramme 14 : Logique du Price Analyzer (Analyse comparative du marché)}
\label{diag:priceanalyzer}
Ce schéma montre comment l'application élabore le badge analytique (ex: *"Moins cher que la moyenne"*) pour aider l'acheteur.

\begin{figure}[H]
\centering
\begin{tikzpicture}[node distance=8mm, scale=0.9, every node/.style={scale=0.9}]

  \node[node-blue] (start) {Consultation de la fiche d'un appartement};
  
  \node[node-white, below=5mm of start, text width=7.5cm] (avg) {
    Récupération du prix moyen dans \texttt{neighborhood\_analytics}\\
    pour le type de bien dans le même quartier
  };

  \node[base-node, shape=diamond, aspect=2.2, below=6mm of avg] (dec) {Prix de l'appartement vs Moyenne ?};

  \node[node-green, below left=1cm and 2mm of dec, text width=3.8cm] (lower) {
    \textbf{Inférieur à la Moyenne}\\
    \scriptsize Badge : "Offre économique"
  };

  \node[node-gold, below=1cm of dec, text width=3.8cm] (mean) {
    \textbf{Égal à la moyenne ($\pm 10\%$)}\\
    \scriptsize Badge : "Au juste prix"
  };

  \node[node-red, below right=1cm and 2mm of dec, text width=3.8cm] (higher) {
    \textbf{Supérieur à la Moyenne}\\
    \scriptsize Badge : "Bâtiment Premium"
  };

  % Liaisons
  \draw[flow-arrow] (start) -- (avg);
  \draw[flow-arrow] (avg) -- (dec);
  \draw[flow-arrow] (dec.south west) -- (lower.north);
  \draw[flow-arrow] (dec.south) -- (mean.north);
  \draw[flow-arrow] (dec.south east) -- (higher.north);

\end{tikzpicture}
\caption{Processus d'analyse de valeur immobilière locale.}
\end{figure}

\newpage

% ================================================================== %
% DIAGRAMME 15 : TRADUCTION DECRIPTIONS (NOUVEAU)
% ================================================================== %
\section*{Diagramme 15 : Pipeline de traduction et d'internationalisation}
\label{diag:traductions}
Comment les annonces sont traduites à la volée entre le Français et l'Anglais.

\begin{figure}[H]
\centering
\begin{tikzpicture}[node distance=8mm, scale=0.9, every node/.style={scale=0.9}]

  \node[node-blue] (request) {Requête HTTP cliente avec en-tête\\ \texttt{Accept-Language: en}};
  
  \node[node-white, below=6mm of request, text width=6.5cm] (api) {FastAPI intercepte la requête};

  \node[base-node, shape=diamond, aspect=2.2, below=6mm of api] (dec) {Traduction présente en base ?};

  \node[node-green, below left=8mm and 5mm of dec, text width=5cm] (db) {
    Lecture de la table \texttt{listing\_translations}\\
    \scriptsize Chargement immédiat de la traduction
  };

  \node[node-gold, below right=8mm and 5mm of dec, text width=5cm] (translate) {
    Appel du traducteur arrière\\
    \scriptsize Traduit le français en anglais, stocke en base
  };

  \node[node-blue, below=4cm of dec, text width=6cm] (response) {Retour de l'annonce en anglais};

  % Liaisons
  \draw[flow-arrow] (request) -- (api);
  \draw[flow-arrow] (api) -- (dec);
  \draw[flow-arrow] (dec.south west) -- (db.north);
  \draw[flow-arrow] (dec.south east) -- (translate.north);
  \draw[flow-arrow] (db.south) -- ++(0,-5mm) -| (response.north);
  \draw[flow-arrow] (translate.south) -- ++(0,-5mm) -| (response.north);

\end{tikzpicture}
\caption{Fonctionnement de la traduction et du bilinguisme dynamique.}
\end{figure}

\end{document}
